%%%%%%%%%%%%%%%%%%%%%%%%%%%%%%%%%%%%%%%%%%%%%%%%%%%%%%%%%%%%%%%%%%%%%%%%%%
% PARTE 1: IMPOSTAZIONI GENERALI DEL DOCUMENTO
%%%%%%%%%%%%%%%%%%%%%%%%%%%%%%%%%%%%%%%%%%%%%%%%%%%%%%%%%%%%%%%%%%%%%%%%%%

\documentclass[11pt, a4paper]{report}

% LINGUA E CODIFICA
\usepackage[italian]{babel}
\usepackage[utf8]{inputenc}
\usepackage[T1]{fontenc}
\usepackage{lmodern}

% LAYOUT DI PAGINA E TIPOGRAFIA
\usepackage[margin=1.5cm, top=1.5cm, bottom=1.5cm]{geometry}
\usepackage{parskip}
\usepackage{microtype}

% GESTIONE DEI COLORI
\usepackage[table]{xcolor}

% MATEMATICA E SIMBOLI
\usepackage{amsmath, amssymb, amsfonts}
\usepackage{mathtools}
\usepackage{siunitx}

% BB OTTIMIZZAZIONE MATRICI LUNGHE PER LA STAMPA
\setlength{\arraycolsep}{10pt}
\usepackage{bm}

% GRAFICA, IMMAGINI E TABELLE
\usepackage{graphicx}
\graphicspath{{./Images/Tesina}}
\usepackage{float}
\usepackage{booktabs}
\usepackage{array}
\usepackage{multirow}
\usepackage{tabularx}

% DIAGRAMMI (TIKZ)
\usepackage{tikz}
\usetikzlibrary{shapes.geometric, arrows.meta, positioning, calc, quotes}
\usepackage{dirtree}

% LINK, RIFERIMENTI ED ELENCHI
\usepackage[unicode]{hyperref}
\hypersetup{
    colorlinks=true,
    linkcolor=black,
    urlcolor=blue,
    citecolor=blue,
    pdftitle={Tesina Industrial Robotic},
    pdfauthor={Francesco Lo Verde},
    pdfpagemode=FullScreen,
}
\usepackage{bookmark}
\usepackage[italian,capitalize]{cleveref}
\usepackage{enumitem}
\setlist{noitemsep}

% PERSONALIZZAZIONE TITOLI, TEOREMI E ALTRI PACCHETTI
\usepackage{titlesec}
\titleformat{\chapter}{\normalfont\huge\bfseries}{\thechapter.}{20pt}{\huge}
\titleformat{\section}{\normalfont\LARGE\bfseries}{\thesection}{1em}{}
\titleformat{\subsection}{\normalfont\Large\bfseries}{\thesubsection}{1em}{}
\titleformat{\subsubsection}{\normalfont\large\bfseries}{\thesubsubsection}{1em}{}
\titlespacing*{\chapter}{0pt}{0pt}{10pt}

\usepackage{comment}

%%%%%%%%%%%%%%%%%%%%%%%%%%%%%%%%%%%%%%%%%%%%%%%%%%%%%%%%%%%%%%%%%%%%%%%%%%
% PARTE 2: DEFINIZIONE DEGLI STILI PER IL CODICE
%%%%%%%%%%%%%%%%%%%%%%%%%%%%%%%%%%%%%%%%%%%%%%%%%%%%%%%%%%%%%%%%%%%%%%%%%%

\usepackage{listings}

% Definizione palette colori (ispirata a VS Code)
\definecolor{codepurple}{RGB}{204, 120, 220}
\definecolor{codecyan}{RGB}{78, 201, 176}
\definecolor{codeorange}{RGB}{206, 145, 120}
\definecolor{codegreen}{RGB}{106, 153, 85}
\definecolor{codeblue}{RGB}{86, 156, 214}
\definecolor{codegray}{rgb}{0.5, 0.5, 0.5}

% Stile di base comune a tutti i linguaggi
\lstdefinestyle{baseStyle}{
    backgroundcolor={},
    commentstyle=\color{codegreen}\itshape,
    stringstyle=\color{codeorange},
    basicstyle=\footnotesize\ttfamily,
    breaklines=true,
    breakatwhitespace=false,
    keepspaces=true,
    numbers=none,
    frame=none,
    showspaces=false,
    showstringspaces=false,
    showtabs=false,
    tabsize=4,
    captionpos=b,
    keywordstyle=\bfseries,
}

% Stile per BASH
\lstdefinestyle{bashStyle}{
    style=baseStyle,
    language=bash,
    keywordstyle={\color{codeblue}\bfseries},
    morekeywords={echo, ls, cd, mkdir, rm, cp, mv, sudo, apt-get, if, then, else, fi, for, in, do, done, while, case, esac, function},
}

% Stile per PYTHON
\lstdefinestyle{pythonStyle}{
    style=baseStyle,
    language=Python,
    keywordstyle=[1]{\color{codepurple}\bfseries},
    keywordstyle=[2]{\color{codeblue}},
    keywordstyle=[3]{\color{codecyan}\itshape},
    morekeywords={
        True, False, None, self, as, with, and, or, not, is, in,
        def, class, import, from, return, if, else, elif, for, while,
        try, except, finally, raise, pass, lambda, yield, global, nonlocal
    },
    morekeywords=[2]{
        print, len, range, int, str, float, list, dict, set, tuple, type,
        open, super, isinstance, zip, enumerate, map, filter
    },
    morekeywords=[3]{
        __init__, __str__, __repr__
    },
    sensitive=true,
}

% Stile per XML
\lstdefinestyle{XMLStyle}{
    style=baseStyle,
    language=XML,
    commentstyle=\color{codegreen},
    keywordstyle=\color{codeblue},
    morekeywords={version, encoding},
    otherkeywords={=},
    stringstyle=\color{codeorange},
    tagstyle=\color{codeblue},
    ndkeywordstyle=\color{codecyan},
}

\author{Francesco Lo Verde}
\begin{document}
% --- Pagina del Titolo ---
\begin{titlepage}
    \centering

    % --- SEZIONE SUPERIORE: LOGO E UNIVERSITÀ ---
    \includegraphics[width=0.35\textwidth]{Images/Tesina/LogoUnipg.png}\par\vspace{1.5cm}

    {\scshape\Large Università degli Studi di Perugia\par}

    \vfill % Aggiunge spazio flessibile

    % --- SEZIONE CENTRALE: TITOLI ---
    % Il titolo è lungo, quindi \huge è una buona dimensione.
    % Usiamo una minipage per controllare meglio la larghezza e prevenire una sillabazione sgradevole.
    \begin{minipage}{0.9\textwidth}
        \centering
        {\Huge \bfseries Sviluppo Di Un Package ROS Per La Modellazione E L'Analisi Cinematica Di Un Manipolatore Di Stanford (6 DOF)\par}
    \end{minipage}

    \vspace{1cm}
    \rule{1\textwidth}{0.4pt} % Linea di separazione
    \vspace{1cm}

    {\LARGE \bfseries Corso di Industrial Robotics\par}

    \vfill % Aggiunge altro spazio flessibile

    % --- SEZIONE INFERIORE: AUTORE ---
    \begin{minipage}{0.7\textwidth}
    \begin{flushleft}
        \large
        \begin{tabular}{@{} l l}
            \multicolumn{2}{l}{\textbf{Studente}} \\
            \addlinespace[0.3em]
            Nome: & \textbf{Francesco Lo Verde} \\
            Email: & \texttt{francesco.loverde@studenti.unipg.it} \\
        \end{tabular}
    \end{flushleft}
    \end{minipage}

    \vspace{1.5cm}

    % --- PIÈ DI PAGINA: ANNO ACCADEMICO ---
    {\large Anno Accademico 2024-2025\par}
    %\vspace{0.5cm}
    %{\small Data ultimo aggiornamento: \today\par}

\end{titlepage}

\restoregeometry

% --- Indice ---
\tableofcontents
\cleardoublepage

%%%%%%%%%%%%%%%%%%%%%%%%%%%%%%%%%%%%%%%%%%%%%%%%%%%%%%%%%%%%%%%%%%%%%%%%%%
% INIZIO CAPITOLO 1
%%%%%%%%%%%%%%%%%%%%%%%%%%%%%%%%%%%%%%%%%%%%%%%%%%%%%%%%%%%%%%%%%%%%%%%%%%

\chapter{Introduzione}
\section{Obiettivi Del Progetto}
Questo report documenta lo sviluppo di un package \textbf{Robot Operating System (ROS)} per la simulazione e l'analisi di un manipolatore robotico a sei gradi di libertà. L'obiettivo primario è la creazione di un ambiente che permetta di validare il modello cinematico del robot e di analizzarne la cinematica diretta. Tale ambiente funge da base per futuri sviluppi, quali la pianificazione del movimento (cinematica inversa) e il controllo avanzato, abilitando il potenziale moto autonomo e preciso del manipolatore.

Gli obiettivi specifici del progetto sono i seguenti:
\begin{itemize}
    \item Sviluppare un modello cinematico e visivo del robot utilizzando il formato \texttt{URDF} (\textit{Unified Robot Description Format}).

    \item Configurare un ambiente di visualizzazione 3D in tempo reale mediante lo strumento \texttt{RViz}, per consentire l'ispezione del modello e il controllo interattivo della configurazione dei giunti.

    \item Realizzare file di lancio (\texttt{.launch}) per l'avvio di tutti i nodi ROS necessari, garantendo che il modello, lo stato dei giunti e l'ambiente di visualizzazione siano caricati e interconnessi correttamente.

    \item Implementare uno script in Python per l'analisi cinematica, in grado di calcolare e visualizzare sul terminale la matrice di trasformazione omogenea \(4 \times 4\), che descrive la posa (posizione e orientamento) dell'end-effector rispetto al sistema di riferimento globale in ogni istante.
\end{itemize}

\subsection{Il Manipolatore Di Stanford}
Il modello del manipolatore sviluppato in questo progetto si ispira all'architettura classica del \textbf{manipolatore di Stanford}. Questa configurazione cinematica a sei gradi di libertà (DOF) costituisce un caso di studio di notevole interesse didattico e analitico.

\begin{figure}[h!]
    \centering
    \begin{minipage}{0.35\textwidth}
        \includegraphics[width=\linewidth]{StanfordManipulator.jpeg}
    \end{minipage}%
    \begin{minipage}{0.65\textwidth}
        \vspace{0pt}
        La sua catena cinematica è composta da:
        \begin{itemize}
            \item Un \textbf{giunto rotoidale} per la rotazione della base (asse verticale).
            \item Un secondo \textbf{giunto rotoidale} per l'elevazione della spalla (asse orizzontale).
            \item Un \textbf{giunto prismatico} per l'estensione del braccio lungo il suo asse principale.
            \item Un \textbf{polso sferico} composto da tre giunti rotoidali, i cui assi si intersecano in un unico punto, che conferisce all'end-effector tre gradi di libertà rotazionali.
        \end{itemize}
    \end{minipage}
\end{figure}

Questa architettura, pur essendo sufficientemente complessa da rappresentare, possiede una cinematica ben nota e analiticamente trattabile.


%%%%%%%%%%%%%%%%%%%%%%%%%%%%%%%%%%%%%%%%%%%%%%%%%%%%%%%%%%%%%%%%%%%%%%%%%%
% INIZIO CAPITOLO 2
%%%%%%%%%%%%%%%%%%%%%%%%%%%%%%%%%%%%%%%%%%%%%%%%%%%%%%%%%%%%%%%%%%%%%%%%%%
\chapter{Strumenti E Architettura Del Sistema}
\label{chap:strumenti_tecnologie}

\section{Il Framework ROS (Robot Operating System)}
L'architettura di \textbf{ROS (Robot Operating System)} è basata su nodi che comunicano tramite paradigmi come topic (publish/subscribe) e service (request/reply), e questo ha permesso di strutturare il progetto in componenti modulari.

La scelta di ROS è motivata dall'elevato numero di strumenti specifici per la robotica che fornisce. L'impiego del visualizzatore \textbf{RViz} e della libreria di gestione delle trasformazioni \textbf{TF2} ha permesso di evitare di implementare la logica di visualizzazione e calcolo, a favore dello studio della cinematica del manipolatore. Il grafo di comunicazione risultante, visualizzabile tramite \texttt{rqt\_graph}, evidenzia l'interazione tra i nodi principali del sistema.
% rqt_graph
\begin{figure}[h!]
   \centering
   \includegraphics[width=\textwidth]{ROS_Graph.png}
   \caption{Diagramma Nodi (Ellissi) E Topic (Rettangoli) Attivi generato con \texttt{rqt\_graph}.}
\end{figure}

\subsection{Modellazione Del Robot: URDF}

L'\textbf{Unified Robot Description Format (URDF)} è lo standard XML adottato per la definizione della topologia del robot. Il modello è descritto come un albero cinematico composto da:
\begin{itemize}
    \item \textbf{Link:} Corpi rigidi definiti dalle loro proprietà inerziali, visive e di collisione.
    \item \textbf{Joint:} Elementi che definiscono la relazione cinematica tra un link genitore (\textit{parent}) e uno figlio (\textit{child}), specificandone il tipo di moto (\texttt{revolute}, \texttt{prismatic}, \texttt{fixed}), l'asse e i limiti.
\end{itemize}

\subsection{Gestione Delle Trasformazioni}
La gestione delle relazioni spaziali dinamiche tra i sistemi di coordinate (\textbf{frame}) è delegata alla libreria \textbf{TF2}. Essa mantiene un albero di trasformazioni (\textbf{TF tree}) aggiornato in tempo reale, che rappresenta la posa di ogni frame del robot rispetto agli altri.

Per quanto riguarda questo progetto:
\begin{itemize}
    \item Il nodo \texttt{robot\_state\_publisher} popola il TF tree a partire dai dati pubblicati sul topic \texttt{/joint\_states} e dalla struttura definita nel modello URDF.

    \item Nodi come RViz e lo script di analisi agiscono invece da \textbf{TF Listener}, interrogando l'albero per ottenere la trasformazione tra due frame qualsiasi in un dato istante.
\end{itemize}

Lo script Python \texttt{tfPrinter.py} implementa un TF Listener il cui scopo è ottenere la trasformazione tra il sistema di riferimento dell'end-effector (\texttt{EndEffector}) e quello globale (\texttt{World}), calcolando la matrice di trasformazione omogenea \(^{World}T_{EndEffector}\).

\begin{figure}[H]
   \centering
   \includegraphics[width=0.5\textwidth]{TF_Tree}
   \caption{Albero delle trasformazioni (TF Tree) del manipolatore.}
   \label{fig:tf_tree}
\end{figure}
\newpage

\section{Struttura e Implementazione del Package \texttt{robot}}
L'intero progetto è incapsulato nel package ROS \texttt{robot}, sviluppato secondo le convenzioni standard. Il package contiene tutti i file di configurazione, i modelli e gli script necessari per l'esecuzione del progetto. La sua organizzazione interna è la seguente:

\dirtree{%
.1 robot.
.2 launch.
.3 robot.launch.
.2 rviz.
.3 robot.rviz.
.2 scripts.
.3 tfPrinter.py.
.2 urdf.
.3 robot.urdf.xacro.
.2 CMakeLists.txt.
.2 package.xml.
}

\subsection{Gestione Del Sistema Tramite File \texttt{launch/}}
\label{subsec:file_di_lancio}
In ROS, i file di lancio (\texttt{.launch}) sono script XML utilizzati per avviare e configurare più nodi contemporaneamente. Il file \texttt{robot.launch} del progetto istanzia l'intero ambiente di simulazione con un singolo comando. La sua esecuzione svolge le seguenti operazioni in sequenza:
\begin{enumerate}
    \item \textbf{Caricamento Del Modello:} Il file XACRO del robot, il cui percorso è definito nell'argomento \texttt{model}, viene processato dal comando \texttt{xacro}. L'output, un file URDF standard, viene caricato sul \textbf{Parameter Server} di ROS sotto il parametro \texttt{robot\_description}. Questo rende la struttura del robot accessibile a tutti i nodi del sistema.

    \item \textbf{Pubblicazione Stato Dei Giunti:} Viene avviato un nodo per la gestione dei giunti. Grazie agli argomenti condizionali \texttt{if} e \texttt{unless}, il file di lancio può avviare due nodi alternativi:
    \begin{itemize}
        \item \texttt{joint\_state\_publisher\_gui}: Avviato di default, crea un'interfaccia grafica con degli slider che permette all'utente di modificare manualmente la posizione di ciascun giunto del robot.
        \item \texttt{joint\_state\_publisher}: Avviato se \texttt{gui="False"}. È una versione non grafica che pubblica semplicemente i valori di default o zero per ogni giunto.
    \end{itemize}
    Entrambi i nodi pubblicano le informazioni sul topic \texttt{/joint\_states}.

    \item \textbf{Calcolo delle Trasformazioni (TF):} Il nodo \texttt{robot\_state\_publisher} si sottoscrive al topic\\
    \texttt{/joint\_states} e utilizzando \texttt{robot\_description}, calcola in tempo reale la posa (posizione e orientamento) di ciascun link rispetto al suo genitore.

    \item \textbf{Avvio di RViz:} Viene lanciato il nodo \texttt{rviz}. L'argomento \texttt{-d \$(find robot)/rviz/robot.rviz} è fondamentale: istruisce RViz a caricare all'avvio una configurazione predefinita, garantendo che la scena sia sempre preparata nel modo corretto.

    \item \textbf{Avvio Script di Analisi:} Viene eseguito il nodo che avvia lo script Python \texttt{tfPrinter.py}. L'opzione \texttt{output="screen"} assicura che i risultati vengano stampati direttamente nel terminale.
\end{enumerate}
\newpage
\subsection{Interfaccia Di Visualizzazione E Controllo (\texttt{rviz/})}
\label{subsec:configurazione_rviz}
L'interazione con il modello del robot avviene principalmente attraverso l'interfaccia di RViz, la cui configurazione è salvata nel file \texttt{robot.rviz}. Questo file memorizza tutti gli elementi della scena, garantendo un'esperienza utente consistente. La configurazione caricata imposta i seguenti elementi:
\begin{itemize}
    \item Imposta il sistema di riferimento fisso su \texttt{World}. Questo definisce l'origine immobile rispetto alla quale tutte le altre pose vengono visualizzate.
    \item Aggiunge un plugin (\texttt{RobotModel}) di visualizzazione che legge la descrizione del robot dal parametro \texttt{robot\_description} e ne visualizza il modello 3D, utilizzando le geometrie e i colori definiti nel file URDF.
    \item Abilita la visualizzazione dei frame per ciascun link del robot. Questo strumento è molto utile per verificare l'orientamento e la posizione di ogni link siano corretti.
\end{itemize}

Il controllo interattivo del manipolatore è reso possibile dall'interazione tra i vari componenti:
\begin{enumerate}
    \item L'utente cambia il valore delle variabili di giunto tramite gli slider.
    \item I nuovi valori vengono pubblicati sul topic \texttt{/joint\_states}.
    \item Il \textbf{\texttt{robot\_state\_publisher}} riceve questi aggiornamenti, ricalcola l'albero TF e lo ripubblica.
    \item \textbf{RViz}, che è sottoscritto al TF tree, aggiorna istantaneamente la posa del modello 3D del robot a schermo.
\end{enumerate}
Questo permette di testare i limiti dei giunti e di validare visivamente la correttezza del modello cinematico prima di procedere con analisi più complesse.

\subsection{Script Di Analisi Cinematica (\texttt{scripts/})}
\label{subsec:scripts}
Il file in questa directory è \texttt{tfPrinter.py}, che agisce come un nodo ROS dedicato all'analisi della cinematica diretta. Implementa un \textbf{TF Listener} che si sottoscrive all'albero delle trasformazioni e richiede periodicamente la posa del frame \texttt{EndEffector} rispetto al frame \texttt{World}. Una volta ottenuti i dati di traslazione e rotazione, calcola la corrispondente matrice di trasformazione omogenea e la stampa a terminale, fornendo così una validazione numerica e in tempo reale del modello cinematico.

\subsection{Definizione del Modello (\texttt{urdf/})}
La directory \texttt{urdf} contiene la definizione geometrica e cinematica del manipolatore. In particolare, il file \texttt{robot.urdf.xacro} descrive il manipolatore Stanford a 6 DOF. Definisce ogni \textbf{link} (con le sue proprietà visive) e ogni \textbf{joint} (con il suo tipo, asse di moto e limiti), costruendo l'intero albero cinematico del robot.

\subsection{Configurazione E Dipendenze (\texttt{package.xml} E \texttt{CMakeLists.txt})}
Alla radice del package si trovano i file standard per la sua gestione.
\begin{itemize}
    \item \texttt{package.xml}: Il file manifesto del package. Dichiara i metadati (nome, versione, autore) e, crucialmente, le dipendenze da altri package ROS (\texttt{rospy}, \texttt{tf2\_ros}, etc.).
    \item \texttt{CMakeLists.txt}: File di build che istruisce il sistema di compilazione Catkin su come costruire il package, specificando quali file (script, launch, modelli) devono essere installati per renderli eseguibili e rintracciabili da ROS.
\end{itemize}


%%%%%%%%%%%%%%%%%%%%%%%%%%%%%%%%%%%%%%%%%%%%%%%%%%%%%%%%%%%%%%%%%%%%%%%%%%
% INIZIO CAPITOLO 3
%%%%%%%%%%%%%%%%%%%%%%%%%%%%%%%%%%%%%%%%%%%%%%%%%%%%%%%%%%%%%%%%%%%%%%%%%%
\chapter{Modellazione Cinematica Del Manipolatore}
\label{chap:modellazione_cinematica}

La modellazione cinematica costituisce il cuore dell'analisi di un manipolatore robotico. Essa permette di descrivere matematicamente la relazione tra i valori dei giunti del robot e la posa (posizione e orientamento) del suo organo terminale (end-effector). Per questo progetto, si è adottato un approccio basato sulla convenzione di Denavit-Hartenberg.

La \textbf{convenzione di Denavit-Hartenberg (DH)} è una metodologia standardizzata per assegnare sistemi di riferimento solidali a ciascun link di una catena cinematica aperta. Il suo vantaggio principale risiede nel descrivere la trasformazione spaziale tra due link consecutivi utilizzando un set minimo di quattro parametri.

Seguendo la procedura DH, si associa un sistema di riferimento \(RF_i\) a ciascun link \(i\). La trasformazione dal sistema di riferimento \(RF_{i-1}\) al sistema \(RF_i\) è definita da quattro parametri:
\begin{itemize}
    \item \(\bm{\theta_i}\) (\textit{Angolo Di Giunto}): La rotazione attorno all'asse \(Z_{i-1}\) per allineare l'asse \(X_{i-1}\) con l'asse \(X_i\).
    \item \(\bm{d_i}\) (\textit{Offset Del Link}): La traslazione lungo l'asse \(Z_{i-1}\) per portare l'origine di \(RF_{i-1}\) sul piano \(X_{i}Y_{i}\).
    \item \(\bm{a_i}\) (\textit{Lunghezza Del Link}): La traslazione lungo il nuovo asse \(X_i\) per portare l'origine di \(RF_{i-1}\) all'origine di \(RF_i\).
    \item \(\bm{\alpha_i}\) (\textit{Torsione Del Link}): La rotazione attorno al nuovo asse \(X_i\) per allineare l'asse \(Z_{i-1}\) con l'asse \(Z_i\).
\end{itemize}
Per giunti rotoidali, \(\theta_i\) è la variabile di giunto, mentre per giunti prismatici lo è \(d_i\). Questi quattro parametri permettono di costruire la matrice di trasformazione omogenea \(^{i-1}T_i\) che descrive la posa del frame \(i\) rispetto al frame \(i-1\).

L'idea fondamentale della convenzione D-H è che la trasformazione spaziale complessa tra due link può essere scomposta in una sequenza standardizzata di quattro movimenti elementari, ciascuno descritto da uno dei parametri. La composizione di questi quattro movimenti (due rotazioni e due traslazioni) attraverso il prodotto delle rispettive matrici porta alla formulazione della seguente matrice di trasformazione omogenea generale:
\[
    ^{i-1}T_i =
    \begin{pmatrix}
        \cos\theta_i & -\sin\theta_i\cos\alpha_i & \sin\theta_i\sin\alpha_i & a_i\cos\theta_i \\
        \sin\theta_i & \cos\theta_i\cos\alpha_i & -\cos\theta_i\sin\alpha_i & a_i\sin\theta_i \\
        0 & \sin\alpha_i & \cos\alpha_i & d_i \\
        0 & 0 & 0 & 1
    \end{pmatrix}
\]

Applicando questa convenzione al manipolatore di Stanford a 6 gradi di libertà, si ottengono i parametri che sono riassunti nella Tabella:

\begin{table}[h!]
    \centering
    \begin{tabular}{|c|c|c|c|c|c|}
        \hline
        \textbf{Trasformazione} & \textbf{\# Link} & \(\bm{\theta_i}\) & \(\bm{d_i}\) & \(\bm{\alpha_i}\) & \(\bm{a_i}\) \\
        \hline
        $RF_{0} \to RF_{1}$ & 1 & $\theta_{1}$ & $d_{1}$ & $-\frac{\pi}{2}$ & 0 \\
        \hline
        $RF_{1} \to RF_{2}$ & 2 & $\theta_{2}$ & $d_{2}$ & $\frac{\pi}{2}$ & 0 \\
        \hline
        $RF_{2} \to RF_{3}$ & 3 & $0$ & $d_{3}$ & 0 & 0 \\
        \hline
        $RF_{3} \to RF_{4}$ & 4 & $\theta_{4}$ & $0$ & $-\frac{\pi}{2}$ & 0\\
        \hline
        $RF_{4} \to RF_{5}$ & 5 & $\theta_{5}$ & $0$ & $\frac{\pi}{2}$ & 0\\
        \hline
        $RF_{5} \to RF_{6}$ & 6 & $\theta_{6}$ & $d_6$ & 0 & 0\\
        \hline
    \end{tabular}
    \caption{Parametri Di Denavit-Hartenberg Per Il Manipolatore Di Stanford A 6 DOF.}
    \label{tab:dh_stanford}
\end{table}
\begin{figure}[H]
   \centering
   \includegraphics[width=0.9\textwidth]{Frame_Chain.png}
   \caption{Assegnazione Dei Frame URDF Secondo La Convenzione DH.}
   \label{fig:dh_frames}
\end{figure}
L'analisi della tabella rivela la struttura cinematica del robot:
\begin{itemize}
    \item \textbf{Giunti 1 e 2:} Corrispondono alla "vita" e alla "spalla" del manipolatore. Sono giunti rotoidali, con variabili \(\theta_1\) e \(\theta_2\). Le torsioni \(\alpha_1 = -\pi/2\) e \(\alpha_2 = \pi/2\) sono necessarie per passare dall'asse di rotazione verticale del primo giunto a quello orizzontale del secondo. I parametri costanti \(d_1\) e \(d_2\) rappresentano gli offset strutturali tra i giunti.
    \item \textbf{Giunto 3:} È il giunto prismatico, responsabile dell'estensione del braccio. La sua variabile di giunto è infatti \(d_3\). L'angolo \(\theta_3\) è nullo, indicando l'assenza di rotazione.
    \item \textbf{Giunti 4, 5 e 6:} Costituiscono il polso sferico. Sono tre giunti rotoidali con variabili \(\theta_4, \theta_5, \theta_6\). La sequenza di torsioni \(\alpha_4 = -\pi/2\) e \(\alpha_5 = \pi/2\) è caratteristica di un polso i cui assi di rotazione si intersecano, permettendo di orientare liberamente l'end-effector. Il parametro \(d_6\) rappresenta la distanza finale dal centro del polso alla punta dell'organo terminale.
\end{itemize}

La cinematica diretta del manipolatore, ovvero la posa dell'end-effector (frame 6) rispetto alla base (frame 0), si ottiene moltiplicando in sequenza le matrici di trasformazione individuali:
\[
    ^{0}T_6 = {^{0}T_1}(\theta_1) \cdot {^{1}T_2}(\theta_2) \cdot {^{2}T_3}(d_3) \cdot {^{3}T_4}(\theta_4) \cdot {^{4}T_5}(\theta_5) \cdot {^{5}T_6}(\theta_6)
\]
Questo prodotto matriciale fornisce la matrice di trasformazione omogenea completa del robot.
\newpage
\section{Modello Nell'Ambiente ROS}
\label{sec:implementazione_ros}

La traduzione del modello matematico di Denavit-Hartenberg in un formato comprensibile da ROS, come l'URDF, non è un processo diretto, in quanto l'URDF non utilizza nativamente i parametri DH. È necessario, pertanto, interpretare la cinematica definita dalla Tabella DH e mapparla sulla struttura a \texttt{link} e \texttt{joint} richiesta dall'URDF.

Il modello URDF è stato costruito nel file \texttt{robot.urdf.xacro} seguendo questi principi:
\begin{itemize}
    \item \textbf{Definizione Link:} Ogni corpo rigido della catena cinematica, dalla base all'end-effector, è stato definito come un tag \texttt{<link>}. All'interno di ciascun link, il tag \texttt{<visual>} è stato utilizzato per associare una forma geometrica semplice e un colore, al fine di creare una rappresentazione 3D riconoscibile in RViz.
    \item \textbf{Definizione Giunti:} La relazione spaziale tra due link consecutivi, descritta da una matrice di trasformazione DH, è stata implementata tramite un tag \texttt{<joint>}. Ogni giunto connette un \texttt{<parent link>} a un \texttt{<child link>}. La posa statica del frame del figlio rispetto a quello del genitore è definita nel tag \texttt{<origin>}, i cui attributi \texttt{xyz} e \texttt{rpy} (roll, pitch, yaw) sono calcolati a partire dai parametri DH fissi. La variabile di giunto (\(\theta_i\) o \(d_i\)) è gestita specificando il tipo di giunto (\texttt{revolute} o \texttt{prismatic}) e il suo asse di rotazione/traslazione nel tag \texttt{<axis>}. Inoltre, sono stati definiti i limiti di movimento (\texttt{<limit>}) per ciascun giunto, specificando gli angoli o le estensioni massime e minime.
\end{itemize}
L'uso di \textbf{XACRO} si è rivelato fondamentale in questa fase. I parametri geometrici, come le lunghezze dei link (\(d_1, d_2, d_6\)), sono stati definiti come proprietà all'inizio del file, rendendo il modello parametrico e facilmente modificabile senza dover alterare le definizioni di ogni singolo link o giunto.

A titolo di esempio, il seguente estratto di codice dal file \texttt{robot.urdf.xacro} mostra come i principi descritti vengono applicati per definire un link e il giunto che lo connette al suo genitore:

\begin{lstlisting}[style=XMLStyle]
<!-- Giunto 2 (Spalla - Revolute) -->
<joint name="Giunto2" type="revolute">
      <parent link="RF0"/>
      <child link="RF1"/>
      <origin xyz="0 0 ${d2}" rpy="-${PI/2} 0 0"/>
      <axis xyz="0 0 1"/>
      <limit effort="100" velocity="1.0" lower="-${PI}" upper="0"/>
</joint>

<!-- Link 2 -->
<link name="RF1">
<visual>
      <origin xyz="0 0 0" rpy="0 0 0"/>
      <geometry>
            <sphere radius="0.1" />
      </geometry>
      <material name="armColor"/>
</visual>
</link>
\end{lstlisting}
\newpage
\subsection{Visualizzazione Della Catena Cinematica}
\label{subsec:catena_cinematica_world}

Una volta definito il modello cinematico, questo viene caricato nell'ambiente ROS, dove la sua struttura viene rappresentata come un albero di sistemi di riferimento (TF tree) gestito da \texttt{tf2}. Tale albero traduce la gerarchia di link e giunti definita nel file \texttt{robot.urdf.xacro} in una catena di trasformazioni spaziali dinamiche.

A differenza di un approccio in cui il robot viene ancorato a un frame \texttt{world} esterno, in questo progetto il frame globale è stato definito direttamente all'interno del file URDF.
La catena cinematica completa, come interpretata da ROS, è quindi la seguente:
\begin{enumerate}
    \item \textbf{Frame \texttt{World}:} È il sistema di riferimento radice dell'intero modello. Definito come un link statico all'interno del file XACRO, funge da origine immobile a cui l'intera struttura del manipolatore è collegata.
    \item \textbf{Catena del Manipolatore:} Il primo giunto mobile, \texttt{Giunto1}, connette il frame \texttt{World} al primo link della catena, \texttt{RF0}. Da qui inizia la sequenza di link e giunti che costituisce il manipolatore, riflettendo la struttura cinematica progettata:
    \[
        \texttt{World} \xrightarrow{\text{Giunto1}} \texttt{RF0} \xrightarrow{\text{Giunto2}} \texttt{RF1} \dots \xrightarrow{\text{Giunto6}} \texttt{RF5}
    \]
    Ogni giunto della catena, da \texttt{Giunto1} a \texttt{Giunto6}, è definito con il suo tipo (\texttt{revolute} o \texttt{prismatic}) e i suoi limiti, permettendo al robot di assumere le configurazioni desiderate.
    \item \textbf{Frame \texttt{EndEffector}:} All'ultimo link della catena cinematica principale, \texttt{RF5}, è collegato il frame finale \texttt{EndEffector} tramite un giunto di tipo fisso, \texttt{GiuntoEE}. Questo frame non rappresenta una parte mobile, ma definisce un punto di interesse cruciale: il \textbf{Tool Center Point (TCP)}, ovvero il punto operativo dell'organo terminale. La posa di questo specifico frame è l'obiettivo finale dell'analisi di cinematica diretta.
\end{enumerate}

Il nodo \texttt{robot\_state\_publisher}, avviato dal file di lancio, svolge il ruolo cruciale di "animare" questa catena. Esso legge i valori angolari e lineari pubblicati sul topic \texttt{/joint\_states} provenienti dagli slide visualizzati da \texttt{joint\_state\_publisher\_gui} e, utilizzando il modello URDF, calcola e pubblica continuamente le trasformazioni TF tra tutti i frame collegati. In questo modo, l'albero TF viene aggiornato in tempo reale, riflettendo la configurazione corrente del manipolatore e rendendo possibile l'interrogazione della posa di qualsiasi frame rispetto a un altro.


%%%%%%%%%%%%%%%%%%%%%%%%%%%%%%%%%%%%%%%%%%%%%%%%%%%%%%%%%%%%%%%%%%%%%%%%%%
% INIZIO CAPITOLO 4
%%%%%%%%%%%%%%%%%%%%%%%%%%%%%%%%%%%%%%%%%%%%%%%%%%%%%%%%%%%%%%%%%%%%%%%%%%
\chapter{Analisi Della Trasformazione Cinematica}
\label{chap:analisi_cinematica}
La validazione numerica del modello cinematico rappresenta il passo conclusivo e fondamentale del processo di modellazione. Essa consiste nel verificare che la rappresentazione digitale del robot, implementata nell'ambiente di simulazione, corrisponda accuratamente al modello matematico. A tal fine, è stato sviluppato uno strumento software in Python che si interfaccia con la libreria TF2 di ROS per estrarre e analizzare la posa dell'end-effector in tempo reale.

\section{Acquisizione Della Posa}
\label{sec:acquisizione_dati}
Il primo passo per la validazione numerica del modello consiste nell'estrarre la posa effettiva dell'end-effector dal sistema ROS in esecuzione. Anziché ricalcolare la cinematica da zero, si sfrutta il sistema \textbf{TF2}, che funge da sorgente per tutte le relazioni spaziali.

Lo strumento per l'acquisizione dei dati è lo script \texttt{tfPrinter.py}. La sua architettura è basata sulle classi fornite dalla libreria \texttt{tf2\_ros}: viene istanziato un buffer (\texttt{tf2\_ros.Buffer}) per memorizzare le trasformazioni ricevute e un listener (\texttt{tf2\_ros.TransformListener}) che si sottoscrive attivamente all'albero TF. Il nodo opera all'interno di un ciclo continuo, la cui frequenza è impostata a 1~Hz tramite \texttt{rospy.Rate(1.0)}, garantendo un monitoraggio costante. La logica è inoltre incapsulata in un blocco \texttt{try...except} per gestire elegantemente eventuali eccezioni, come la non disponibilità della trasformazione all'avvio del sistema.

All'interno di ogni iterazione del ciclo, l'operazione fondamentale è la chiamata alla funzione \texttt{lookup\_transform()}. Questa funzione interroga il buffer TF per ottenere la trasformazione più recente disponibile (\texttt{rospy.Time(0)}) tra il sistema di riferimento di destinazione, \texttt{TARGET\_FRAME = 'World'}, e quello di origine, \texttt{SOURCE\_FRAME = 'EndEffector'}.

È fondamentale sottolineare che l'oggetto restituito da questa funzione non è una matrice 4x4, ma una struttura dati che rappresenta la posa attraverso due componenti distinti:
\begin{itemize}
    \item \textbf{Vettore di Traslazione}: Contiene le coordinate \((x, y, z)\) dell'origine del frame \texttt{EndEffector} rispetto all'origine del frame \texttt{World}.
    \item \textbf{Quaternione Unitario}: Descrive l'orientamento del frame \texttt{EndEffector} rispetto al frame \texttt{World} tramite i suoi quattro componenti \((q_x, q_y, q_z, q_w)\).
\end{itemize}

I dati ottenuti da TF2, sebbene efficienti per la comunicazione interna di ROS, non sono nella forma standard utilizzata per l'analisi cinematica. Il passo successivo consiste quindi nel convertire il vettore di traslazione e il quaternione nella rappresentazione matematica più interpretabile: la matrice di trasformazione omogenea.

\subsection{Conversione In Matrice 4x4}
\label{subsec:conversione_matrice}
Nello script \texttt{tfPrinter.py}, il processo di conversione è delegato alla libreria \texttt{tf.transformations}, un modulo Python ampiamente utilizzato nell'ecosistema ROS. La procedura si articola in due fasi:
\begin{enumerate}
    \item \textbf{Conversione del Quaternione:} Il quaternione di rotazione, ottenuto dal listener TF2, viene passato alla funzione \texttt{quaternion\_matrix()}. Questa funzione lo converte nella sua corrispondente matrice di rotazione ortonormale, indicata con \(R\).

    \item \textbf{Assemblaggio della Matrice Finale:} La matrice di trasformazione omogenea 4x4, \(^{World}T_{EndEffector}\), viene assemblata combinando la sottomatrice di rotazione \(R\) (3x3) e il vettore di traslazione (3x1) in un'unica struttura, secondo la forma canonica:
    \[
        ^{World}T_{EndEffector} =
        \begin{pmatrix}
        r_{11} & r_{12} & r_{13} & p_x \\
        r_{21} & r_{22} & r_{23} & p_y \\
        r_{31} & r_{32} & r_{33} & p_z \\
        0 & 0 & 0 & 1
        \end{pmatrix}
    \]
\end{enumerate}
Questa matrice finale incapsula in modo compatto e standardizzato l'intera informazione sulla posa dell'end-effector.

\newpage
La matrice di trasformazione omogenea \(^{World}T_{EndEffector}\) non è un semplice insieme di numeri, ogni sua parte possiede un'interpretazione fisica precisa, fondamentale per comprendere lo stato del robot.
\begin{itemize}
    \item \textbf{Matrice Di Rotazione:} La sottomatrice \(R\) descrive l'\textbf{orientamento} del sistema di riferimento dell'end-effector (\(RF_{EE}\)) rispetto a quello globale (\(RF_{World}\)). In particolare, le sue tre colonne rappresentano i versori degli assi del frame \texttt{EndEffector} espressi nelle coordinate del frame \texttt{World}.
    \[
        R = \begin{pmatrix} \hat{x}_{EE,w} & \hat{y}_{EE,w} & \hat{z}_{EE,w} \end{pmatrix}
    \]
    L'analisi di questa sottomatrice permette di determinare come sono orientati gli assi dell'end-effector.

    \item \textbf{Vettore Di Traslazione:} Il vettore \([p_x, p_y, p_z]^T\) nell'ultima colonna descrive la \textbf{posizione} dell'origine del frame dell'end-effector rispetto all'origine del frame globale. Esso fornisce le coordinate cartesiane del \textit{Tool Center Point} (TCP), specificando dove si trova l'end-effector nello spazio.
\end{itemize}
La combinazione di queste due componenti fornisce una descrizione completa e non ambigua della posa dell'organo terminale, che costituisce il risultato della cinematica diretta.

\section{Validazione Posa Di Riferimento}
\label{sec:validazione_posa}
Per validare la correttezza del modello cinematico e l'accuratezza dello strumento di analisi, si esamina la configurazione più semplice e predicibile del manipolatore: la sua posa di riferimento. Questa configurazione corrisponde allo stato in cui tutti i giunti rotoidali sono a un angolo di zero radianti e il giunto prismatico ha un'estensione nulla. In questo stato, la posa dell'end-effector può essere calcolata manualmente e confrontata con l'output dello script.

Avviando il sistema con il file di lancio \texttt{robot.launch} e senza muovere gli slider, il manipolatore si trova nella sua posa di riferimento. In questa condizione, lo script \texttt{tfPrinter.py} produce il seguente output a terminale:

\begin{verbatim}
-----------------------------------------------------
Trasformazione da 'EndEffector' a 'World'

Timestamp: ...

Traslazione (x, y, z):
     x: 2.2000
     y: 0.6050
     z: 1.9500

Quaternione (Rotazione x, y, z, w) :
     x: 0.0000
     y: 0.7071
     z: 0.0000
     w: 0.7071

Matrice Trasformazione Omogenea World - End Effector:

[[-0.     0.     1.     2.2  ]
 [ 0.     1.     0.     0.605]
 [-1.     0.    -0.     1.95 ]
 [ 0.     0.     0.     1.   ]]
-----------------------------------------------------
\end{verbatim}

\newpage
\subsection{Correlazione Con Il Modello Fisico}
\label{subsec:discussione_risultati}
L'output numerico può essere direttamente correlato alla geometria fisica del robot definita nel file \texttt{robot.urdf.xacro}. L'analisi si scompone nella verifica della posizione e dell'orientamento.

\subsubsection{Analisi del Vettore di Traslazione}
\label{subsec:analisi_vettore_traslazione}
Il vettore di posizione calcolato \(\textbf{p} = [2.2, 0.605, 1.95]^T\) è il risultato della catena di trasformazioni definita nel file URDF. La validazione di questi valori può essere ottenuta tracciando le trasformazioni lungo la catena cinematica:

\begin{enumerate}
    \item Fino al giunto prismatico, il braccio si estende lungo gli assi del frame "World". La posizione dell'origine del frame \texttt{RF2} è la somma degli offset iniziali: una traslazione di \texttt{2.0} metri in X e di \texttt{1.3} metri in Z. Inoltre, l'offset del giunto prismatico aggiunge \texttt{0.2} metri lungo l'asse X. Fino a questo punto, la posizione è \([2.2, 0, 1.3]\)

    \item Il giunto successivo aggiunge un offset verticale di \texttt{0.65} metri. Poiché il frame \texttt{RF2} è ancora allineato con il "World", questa traslazione porta l'origine di \texttt{RF3}  alla posizione
    \[
        \textbf{p} = [2.2, 0, 1.3] + [0, 0, 0.65] = [2.2, 0, 1.95]
    \]

    \item Il \texttt{Giunto5} non trasla l'origine ma introduce una rotazione di 90° attorno all'asse Y locale, che in questo caso corrisponde all'asse Y del frame "World". Da questo momento ogni offset verticale nel frame locale viene visto come un un offset orizzontale nel frame "World".

    \item L'ultimo giunto (\texttt{GiuntoEE}) aggiunge un offset di \texttt{d6 = 0.605} metri lungo l'asse Z del suo frame genitore. Però come conseguenza delle rotazioni accumulate, l'asse Z di \texttt{RF5} non è più verticale ma punta lungo la direzione dell'asse Y del frame "World". Quindi, questo offset sarà una traslazione di \texttt{0.605} metri lungo l'asse Y di quest'ultimo.
\end{enumerate}

Sommando i contributi, si ottiene la posizione finale dell'end-effector:
\[
    \textbf{p} = [2.2, 0, 1.95] + [0, 0.605, 0] = [2.2, 0.605, 1.95]
\]
Avendo trovato gli stessi valori ottenuti dallo script, è dimostrata la coerenza tra il modello URDF e l'output numerico.

\subsubsection{Analisi della Matrice di Rotazione}
La sottomatrice di rotazione \(R\) risultante è:
\[
    R =
    \begin{pmatrix}
        0 & 0 & 1 \\
        0 & 1 & 0 \\
       -1 & 0 & 0
    \end{pmatrix}
\]
Il quaternione \(\mathbf{q} = [0, 0.7071, 0, 0.7071]\) corrisponde a una rotazione di \(\pi/2\) radianti attorno all'asse Y del sistema di riferimento genitore. Questa rotazione è il risultato della sequenza di rotazioni definite nei giunti del robot.

Interpretando fisicamente le colonne della matrice \(R\), si ottiene l'orientamento dell'end-effector rispetto al mondo:
\begin{itemize}
    \item \textbf{Asse X Dell'End-Effector} (Colonna \([0, 0, -1]^T\)): Punta nella direzione opposta all'asse \textbf{Z del Mondo}.
    \item \textbf{Asse Y Dell'End-Effector} (Colonna \([0, 1, 0]^T\)): Punta nella stessa direzione dell'asse \textbf{Y del Mondo}.
    \item \textbf{Asse Z Dell'End-Effector} (Colonna \([1, 0, 0]^T\)): Punta nella stessa direzione dell'asse \textbf{X del Mondo}.
\end{itemize}

%%%%%%%%%%%%%%%%%%%%%%%%%%%%%%%%%%%%%%%%%%%%%%%%%%%%%%%%%%%%%%%%%%%%%%%%%%
% INIZIO CAPITOLO 5
%%%%%%%%%%%%%%%%%%%%%%%%%%%%%%%%%%%%%%%%%%%%%%%%%%%%%%%%%%%%%%%%%%%%%%%%%%
\chapter{Conclusioni}
\label{chap:conclusioni}

\section{Sintesi dei Risultati e Prospettive Future}
\label{sec:sintesi_e_prospettive}

Il progetto ha raggiunto con successo l'obiettivo di realizzare un ambiente di simulazione e analisi completo per un manipolatore robotico a sei gradi di libertà. A tal fine, è stato sviluppato un package ROS che incapsula tutti gli elementi necessari per questo scopo:
\begin{itemize}
    \item Un \textbf{modello cinematico} del robot, definito in formato URDF e reso parametrico e modulare attraverso l'impiego di macro XACRO.
    \item Un \textbf{ambiente di visualizzazione}, orchestrato da un file di lancio, che integra RViz per l'ispezione 3D e un'interfaccia grafica per il controllo manuale dei giunti.
    \item Uno \textbf{strumento di analisi numerica}, implementato come script Python, che si interfaccia con la libreria TF2 per estrarre la posa dell'end-effector in tempo reale.
\end{itemize}

Il risultato più significativo di questo lavoro risiede nella validazione numerica del sistema. L'analisi della posa di riferimento ha confermato una perfetta corrispondenza tra l'output dello script e la geometria definita nel modello. Ciò attesta l'accuratezza della catena cinematica implementata e l'affidabilità dello strumento di verifica.

In conclusione, il sistema sviluppato costituisce una solida base di partenza, che è prerequisito fondamentale per affrontare tematiche più complesse. Gli sviluppi futuri potranno ora spostare il focus dalla pura cinematica al controllo dinamico, attraverso l'integrazione con simulatori fisici come \texttt{Gazebo}, e alla pianificazione intelligente del movimento, mediante l'impiego di framework avanzati come \texttt{MoveIt!}.

\appendix

\chapter{Appendice}
\section{Guida Per L'Avvio}
\label{sec:guida_avvio}

Prima di procedere, è necessario assicurarsi che il sistema soddisfi i seguenti requisiti:
\begin{itemize}
    \item Una distribuzione Linux compatibile con ROS (Ubuntu 20.04).
    \item Una installazione funzionante di \textbf{ROS Noetic}.
    \item Un \textbf{Catkin Workspace} inizializzato. È possibile crearlo con i comandi:
    \begin{lstlisting}[style=bashStyle]
mkdir -p ~/catkin_ws/src
cd ~/catkin_ws/
catkin_make
    \end{lstlisting}
\end{itemize}

A questo punto, la procedura per installare il package \texttt{robot} e le sue dipendenze è:

\begin{enumerate}
    \item Assicurarsi che la cartella del package, denominata \texttt{robot}, sia stata copiata all'interno della directory \texttt{src} del proprio Catkin Workspace.

    \item Compilare l'intero workspace utilizzando il comando \texttt{catkin\_make}. Questo renderà il package e i suoi script disponibili per l'ambiente ROS.
    \begin{lstlisting}[style=bashStyle]
cd ~/catkin_ws
catkin_make
    \end{lstlisting}

    \item Per rendere il package visibile al sistema ROS, è necessario eseguire il "sourcing" del file di setup. È buona norma aggiungerlo al proprio file \texttt{.bashrc} per renderlo permanente.
    \begin{lstlisting}[style=bashStyle]
source ~/catkin_ws/devel/setup.bash
    \end{lstlisting}
\end{enumerate}

Una volta completata l'installazione, l'intero ambiente di simulazione può essere avviato con un singolo comando. Per avviare la simulazione completa, inclusa l'interfaccia grafica per il controllo dei giunti, eseguire il seguente comando:
\begin{lstlisting}[style=bashStyle]
roslaunch robot robot.launch
\end{lstlisting}
Dopo l'esecuzione, dovrebbero aprirsi due finestre:
\begin{itemize}
    \item La finestra principale di \textbf{RViz}, che mostra il modello 3D del manipolatore nella sua posa di riferimento.
    \item Una finestra più piccola intitolata \textbf{"Joint State Publisher"}, contenente gli slider per controllare manualmente ogni giunto.
\end{itemize}
Contemporaneamente, il terminale in cui è stato eseguito il comando inizierà a stampare l'output dello script \texttt{tfPrinter.py}, mostrando la matrice di trasformazione omogenea aggiornata a ogni secondo.

Per avviare il sistema senza la finestra degli slider, è possibile passare un argomento al file di lancio:
\begin{lstlisting}[style=bashStyle]
roslaunch robot robot.launch gui:=False
\end{lstlisting}
In questa modalità, il robot rimarrà fermo nella sua posa di riferimento e lo script di analisi stamperà sempre la stessa matrice di trasformazione.

\newpage
\section{Script Python (tfPrinter.py)}
\label{sec:script_python}
\begin{lstlisting}[style=pythonStyle]
#!/usr/bin/env python3
# -*- coding: utf-8 -*-
import rospy
import tf2_ros
import numpy as np
import tf.transformations as tft
# Impostazioni di stampa per NumPy
# Precisione a 4 decimali, senza notazione scientifica e larghezza di linea
np.set_printoptions(precision=4, suppress=True, linewidth=120)

# Inizializzazione
rospy.init_node('tfPrinter')
# Definizione frame
TARGET_FRAME = 'World'
SOURCE_FRAME = 'EndEffector'
# Creazione buffer e listener TF
tf_buffer = tf2_ros.Buffer()
listener = tf2_ros.TransformListener(tf_buffer)
# Frequenza Di Aggiornamento: 1 volta al secondo
rate = rospy.Rate(1.0)

# Il ciclo principale del nodo ROS
while not rospy.is_shutdown():
      try:
            # rospy.Time(0) per ottenere l'ultima trasformazione disponibile.
            dataTF = tf_buffer.lookup_transform(TARGET_FRAME, SOURCE_FRAME, rospy.Time(0))
            translation = dataTF.transform.translation
            rotation = dataTF.transform.rotation
            # Stampa dei dati grezzi
            print("-----------------------------------------------------")
            print(f"Trasformazione da '{SOURCE_FRAME}' a '{TARGET_FRAME}'")
            print(f"\nTimestamp: {dataTF.header.stamp.to_sec()}")
            print("\nTraslazione (x, y, z):")
            print(f"\t x: {translation.x:.4f}")
            print(f"\t y: {translation.y:.4f}")
            print(f"\t z: {translation.z:.4f}")
            traslationVector = np.array([translation.x, translation.y, translation.z])
            print("\nQuaternione (Rotazione x, y, z, w) :")
            print(f"\t x: {rotation.x:.4f}")
            print(f"\t y: {rotation.y:.4f}")
            print(f"\t z: {rotation.z:.4f}")
            print(f"\t w: {rotation.w:.4f}")
            # Conversione del quaternione in un array NumPy
            quaternionVector = np.array([rotation.x, rotation.y, rotation.z, rotation.w])
            # Calcolo matrice omogernea
            quatMatrix = tft.quaternion_matrix(quaternionVector)
            rotationMatrix = quatMatrix[:3, :3]
            homogeneousMatrix = np.identity(4)
            homogeneousMatrix[:3, :3] = rotationMatrix
            homogeneousMatrix[:3, 3] = traslationVector
            print("\nMatrice Trasformazione Omogenea World - End Effector:\n")
            print(np.round(homogeneousMatrix, 4))

      except (tf2_ros.LookupException, tf2_ros.ConnectivityException, tf2_ros.ExtrapolationException) as e:
            rospy.logwarn_throttle(5.0, "In attesa che la trasformazione diventi disponibile...")

      rate.sleep()
\end{lstlisting}


\end{document}
